\documentclass[10pt,a4paper]{amsart}
\usepackage[margin=19mm]{geometry}
\usepackage{amsmath,amssymb,mathtools}
\usepackage[scr=rsfs,cal=euler]{mathalfa}
\usepackage{xfrac}
\usepackage[unicode,hidelinks,bookmarks=false]{hyperref}
\newcommand{\ie}{i.e.\ }
\newcommand{\cf}{cf.\ }
\newcommand{\resp}{respectively\ }
\newcommand{\Subsection}{\subsection}
\providecommand{\nobreakdash}{\nobreak\hbox{-}}
\setlength{\emergencystretch}{2em}
\multlinegap0pt
\usepackage[shortlabels]{enumitem}
\usepackage{mathtools}
\usepackage{amssymb}
\newtheorem{lemma}{Lemma}
\newcommand{\restr}[2]{#1\big|_{#2}}
\title{From 12 to 6: Sharpening the Three-Charge Bound in Maxwell's Problem}
\author{Andrei Gabrielov, Dmitry Novikov, Tomer Novikov, Boris Shapiro}
\date{Source: arXiv:2607.28785v1 (CC BY 4.0)}
\begin{document}
\maketitle

\noindent\emph{Source and licence.} From 12 to 6: Sharpening the Three-Charge Bound in Maxwell's Problem. Version arXiv:2607.28785v1 (CC BY 4.0), distributed by arXiv under the Creative Commons Attribution 4.0 International licence, http://creativecommons.org/licenses/by/4.0/, as recorded in the arXiv licence metadata for this exact version and shown on the abstract page https://arxiv.org/abs/2607.28785v1. Full licence notice: \texttt{licenses/arxiv-2607.28785-CC-BY-4.0.txt}.\par\smallskip

\noindent\textit{Editorial scope.} Counted unit: the article's lemma lem:four-contacts (source0.tex lines 488--490). The article's complete original proof of it is delivered with it (source0.tex lines 492--537, one \texttt{proof} environment of 46 lines). No mathematics is written, repaired or re-derived: the statement and the proof are the article's own lines, in the article's own notation, and no retained line is substituted or re-flowed.  No further source-local context is needed: every cross-reference of every retained line resolves inside the two delivered spans.  Nothing was omitted from any retained span, so the package carries no omission record.  No retained line contains a citation, so the wrapper carries no bibliography and no bibliography entry.  No retained line contains a figure environment, an image-inclusion command or drawing code, so no image is delivered and the package declares no asset dependency.  Editorial formatting only, in addition: an AMS \texttt{amsart} preamble that restates the article's own declarations and macros used by the retained lines, namely the environments \texttt{lemma} on separate counters, together with the standard packages those lines need.  The retained environments print in this document as Lemma 1 in order of appearance, whereas the article prints the same items with the numbering of its own full document.  The complete native source of the article as distributed in the arXiv e-print for this exact version is bundled unchanged as \texttt{sources/206438/source0.tex}.
\begin{lemma}[Four-contact lemma]\label{lem:four-contacts}
The total multiplicity of the zeros of $R$ on $O$ is at least $4$.
\end{lemma}

\begin{proof}
The function $h=\restr{\Phi}{O}$ is a nonconstant real-analytic function
on a circle.  Its critical points are isolated.  A zero of odd order of
$h'$ is exactly a point at which $h'$ changes sign; hence the number of
odd-order zeros of $h'$ around the circle is even.  Moreover, $h$ attains
its minimum and maximum at two distinct critical points.

Suppose, for a contradiction, that the total multiplicity of the zeros
of $h'$ is at most $3$.  Since $p$ is one of these zeros and there are at
least two distinct critical points, only the following possibilities can
occur:
\begin{enumerate}[label=\textup{(\roman*)},leftmargin=2.4em]
  \item exactly two critical points, both simple;
  \item exactly two critical points, of total multiplicity $3$;
  \item exactly three critical points, all simple.
\end{enumerate}
Cases \textup{(ii)} and \textup{(iii)} are impossible because each gives
an odd number of odd-order zeros of $h'$.

It remains to exclude case \textup{(i)}.  The two critical points are the
unique global minimum and maximum of $h$.  First suppose that $p$ is the
minimum.  Then
\begin{equation}\label{eq:Phi-above}
  \Phi>c_0\qquad\text{on }O\setminus\{p\}.
\end{equation}
The punctured circle $O\setminus\{p\}$ is connected, so the separation
lemma implies that it is contained entirely in $L$ or entirely in
$\mathcal R$.

Because $p$ is a simple critical point of $h$ and a local minimum,
\begin{equation}\label{eq:positive-tangent}
  v_0^{\mathsf T}Hv_0>0.
\end{equation}
The positive cone of the quadratic form
$H=\operatorname{diag}(+,-)$ contains no vertical vector.  Thus the
$f$-component of $v_0$ is nonzero.  The two ends of the arc
$O\setminus\{p\}$ approach $p$ with opposite tangent directions $v_0$
and $-v_0$, and therefore with opposite signs of $f-f_0$.  One end lies
in $L$ and the other in $\mathcal R$, a contradiction.

If $p$ is the maximum, the same argument uses the two components of
$\{\Phi<c_0\}$ separated by the line $g=g_0$; in this case the negative
cone of $H$ contains no horizontal vector.  Hence case
\textup{(i)} is also impossible, and the total multiplicity is at least
$4$.
\end{proof}

\end{document}
